\documentclass[11pt]{article}

\usepackage[margin=1in]{geometry}
\usepackage{amsmath,amssymb,amsthm}
\usepackage{booktabs}
\usepackage{graphicx}
\usepackage{xcolor}
\usepackage{multirow}
\usepackage{hyperref}
\usepackage[numbers,sort&compress]{natbib}
\usepackage{setspace}
\doublespacing
\usepackage{lineno}
\linenumbers

\newtheorem{proposition}{Proposition}

\renewcommand{\arraystretch}{1.3}
\setlength{\textfloatsep}{12pt plus 2pt minus 2pt}
\setlength{\intextsep}{10pt plus 2pt minus 2pt}
\setlength{\abovecaptionskip}{8pt}
\setlength{\belowcaptionskip}{4pt}

\title{Site-level rates and site-level coefficients estimate different\\
quantities: an individual-level benchmark for federated analysis\\
of COVID-19 surveillance data}

\author{Elliot Tower, M.S.\\
\textit{Independent Researcher}\\[6pt]
\textbf{Corresponding author:}\\
Elliot Tower\\
Email: \texttt{elliot@elliottower.ai}}

\date{}

\begin{document}
\maketitle

\begin{abstract}
\noindent\textbf{Objective:}
Federated networks share either site-level rates or coefficients from
site-level models, and treat these as interchangeable.  We tested whether
they estimate the same quantity.

\noindent\textbf{Materials and Methods:}
We analyzed 3\,993\,464 confirmed COVID-19 case records from Mexico's
national surveillance system in 32 treating-unit states.  The observed
ecological slope regressed state case fatality on the share aged 70 or
older.  The model-implied slope averaged individual-level logistic
predictions within each state and regressed those on the same share;
their difference is the ecological discrepancy.  We varied the
individual-level model across six prespecified specifications of
increasing adjustment, and the site definition across residence state,
health-care sector, municipality, and a constrained random partition
holding state size and elderly count fixed.

\noindent\textbf{Results:}
Under the prespecified three-covariate model the observed and
model-implied slopes were 1.31 and 0.45 percentage points of case
fatality per point of elderly share (difference 0.86; bootstrap 95\% CI
0.50--1.30).  The difference fell as adjustment increased but remained
0.59 (0.21--1.07) under the richest specification.  On matched records it
was smaller for municipalities than for states of residence (0.47 vs
1.01), and it was 0.00 under the random partition.  One-stage and
two-stage estimators agreed to 0.0007 in the log odds ratio when their
specifications were aligned.

\noindent\textbf{Discussion:}
The discrepancy is defined relative to a specified individual-level model
and does not identify a contextual effect.

\noindent\textbf{Conclusion:}
Site rates and site coefficients estimate different quantities.  Networks
sharing rates should name the individual-level model their inference
assumes, and share per-site coefficients when an individual-level
association is the target.
\end{abstract}

\noindent\textbf{Keywords:} ecological fallacy; multicenter studies as topic; data aggregation;
electronic health records; COVID-19


\section{Background and Significance}

Federated analysis networks have become central infrastructure for
multi-site clinical research.  The ENACT/ACT network connects
institutions across the CTSA consortium~\cite{visweswaran2018act},
TriNetX provides access to electronic health records from more than
20 countries~\cite{trinetx2024}, and PCORnet integrates clinical
data across multiple partner networks serving over 47 million
patients~\cite{fleurence2014pcornet}.  These
platforms share summary statistics rather than individual records,
and a central analysis pools the summaries into a single estimate.
This approach works when the hospitals are measuring the same thing
in the same kind of patient, but participating sites typically differ
in patient demographics, coding practices, and case severity.
One site serves a young urban population while another draws from
a rural elderly community; one codes aggressively for neurological
complications while another barely records them.  The pooled average
across such sites is a well-defined number, but it may not describe
any actual site or any actual patient---and it can point in the wrong
direction entirely.

This problem---drawing conclusions about individuals from group-level
data---is the \emph{ecological fallacy}
\cite{robinson1950ecological}.  The structural mechanisms behind it---aggregation
bias, cross-level confounding, and specification error---are well understood
\cite{greenland1992ecological,greenland2001ecologic}, and formal frameworks
for ecological inference exist \cite{king1997ecological,wakefield2003ecological,wakefield2008ecologic}.
Known since 1950 and well-covered in epidemiology textbooks, the
ecological fallacy is widely acknowledged yet rarely quantified in
large-scale federated analyses.
Multi-site consortium studies and federated analysis networks typically have
access only to site-level summaries and treat them as proxies for
patient-level truth.

We quantify this distortion using 42 million individual COVID-19 patient
records from two countries, plus a third dataset of 22,000 patients from 21
hospitals across 6 countries.  The CDC dataset uses race/ethnicity categories
as pseudo-sites (7 groups, excluding categories with ambiguous demographic
composition).  Mexico's 32 states provide a confirmatory test on real
geographic sites.  In both cases, ecological regression detects a
statistically significant relationship between age and mortality---but the
ecological coefficient is a linear slope, while the individual-level truth
is a multiplicative odds ratio an order of magnitude larger:

\begin{itemize}
\item \textbf{CDC (38 million records, 7 pseudo-sites):} ecological
$\beta = +0.55$, $p = 0.02$; individual OR $= 9.9$, $p < 0.001$.

\item \textbf{Mexico (4 million records, 32 states):} ecological
$\beta = +1.31$, $p < 0.0001$; individual OR $= 11.3$, $p < 0.001$.
\end{itemize}

Structural consistency tests---methods
that examine the \emph{structure} of relationships across sites, rather than
averaging over them---can diagnose when site-level pooling is unreliable.
Applied to the 4CE consortium's neurological COVID-19 dataset, these methods
detect coordinated heterogeneity invisible to standard meta-analysis.
Federated analyses should treat their pooled summaries as hypotheses for
individual-level testing.


\section{Ecological Bias in Two Datasets}

\subsection{CDC, United States: ecological distortion with pseudo-sites}

The CDC Case Surveillance dataset contains 38.2 million confirmed COVID-19
cases with age group, sex, race/ethnicity, hospitalization status, ICU
admission, and death.  We used 7 race/ethnicity categories as pseudo-sites
(excluding ``Unknown'' and ``Missing'' categories with ambiguous demographic
composition; groups range from $n = 13{,}548$ to $n = 6{,}165{,}036$),
mimicking what a consortium study would have if each racial/ethnic group
were a separate hospital reporting summary statistics.  This is a deliberate
simplification that maximizes compositional confounding; the specific
confounders differ from those in hospital-based consortia (see Limitations).

\textbf{Site-level analysis} aggregated each group's mortality rate
(deaths/cases) and proportion of patients aged 80+, then regressed one on the
other across the 7 groups.  Result: a significant positive relationship
($\beta = +0.55$ mortality-rate units per unit proportion elderly,
$p = 0.02$, $R^2 = 0.69$).  The ecological regression detects that age
predicts mortality.

\textbf{Individual-level analysis} on the same 38.2 million patients found
that being over 80 multiplies the odds of death by 9.9 ($p < 0.001$).
After adjusting for sex and comorbidities, patients gain 2.4 times the odds of
dying for each additional decade of age ($p < 0.001$).  Males have 1.7
times the odds; patients with underlying conditions have 1.5 times the odds.

The ecological regression gets the direction right, but the quantities are
incommensurable (Figure~\ref{fig:ecological}A): a linear slope of $+0.55$
stands in for a multiplicative odds ratio of 9.9.  The ecological
coefficient is not a noisy estimate of the individual-level effect---it is a
different quantity entirely.

Three pre-registered supplementary analyses illuminate the role of the
grouping strategy.  When all 105 million CDC records (including probable
cases and all race/ethnicity categories, compared to the 38.2 million
confirmed cases used in the primary analysis) are randomly
partitioned into 9 groups (destroying the race/ethnicity structure),
ecological regression achieves only 4.8\% power (95\% Wilson CI
[3.2\%, 7.0\%])---indistinguishable from the false-positive rate
(Supplement~S6; Figure~\ref{fig:power}).  The race/ethnicity grouping
succeeds because different racial/ethnic groups have different age
distributions; random grouping averages out these differences, leaving no
signal for ecological regression to detect.  Duncan--Davis ecological
inference bounds on the elderly mortality rate span [0, 24.4\%], containing
the true 9.91\% rate but too wide to be informative (Supplement~S4).


\subsection{Mexico: ecological distortion with real geographic sites}

Mexico's COVID-19 surveillance data provides a confirmatory test on real
geographic sites: 3,993,464 confirmed cases across
32 real geographic states, each corresponding to the state of the treating
medical unit.

\textbf{Site-level analysis} across the 32 states finds a significant
relationship between the proportion of patients aged 70+ and the state
mortality rate ($\beta = +1.31$ mortality-rate units per unit proportion
elderly, $p < 0.0001$, $R^2 = 0.50$).  Concretely, a 10-percentage-point
increase in a state's elderly share predicts a 13.1-percentage-point
increase in state mortality.

\textbf{Individual-level analysis} on the same 4 million patients reveals the
true effect: patients aged 70 or older have 11.3 times the odds of dying
($p < 0.001$).  After adjusting for sex and 9 individual comorbidities
(diabetes, hypertension, obesity, cardiovascular disease, chronic kidney
disease, COPD, asthma, immunosuppression, and smoking), the adjusted odds
ratio remains 8.0 ($p < 0.001$).  Males have 1.7 times the odds;
patients with any comorbidity have 3.8 times the odds.

Site-level analysis gets the direction right and the size wrong (Figure~\ref{fig:ecological}B): a linear
slope of $+1.31$ standing in for an odds ratio of 11.

Supplementary analyses quantify this distortion directly.  To isolate the
aggregation bias, we fitted the individual-level logistic model
(elderly, sex, comorbidity) and computed the ecological slope that
\emph{would} arise if the ecological regression faithfully summarized the
individual-level relationship.  The observed ecological slope is 2.09 times
the expected slope---a 109\% amplification from aggregation bias
(Supplement~S3).  Two individual-level pooling methods---two-stage IPD
meta-analysis (per-state logistic regression pooled via DerSimonian--Laird
random effects) and one-stage IPD (logistic with state fixed
effects)---recover odds ratios of 5.34 and 5.22, respectively, agreeing
within 2\% (Supplement~S9).  Both IPD methods recover the individual-level effect; the ecological regression does not.

\begin{table}[t]
\centering
\caption{Ecological distortion in both datasets.  Site-level regression
detects a significant relationship in both cases, but the ecological
coefficient (a linear slope) is incommensurable with the individual-level
truth (a multiplicative odds ratio).}
\label{tab:ecological}
\begin{tabular}{llcc}
\toprule
Dataset & Analysis level & Effect size & $p$-value \\
\midrule
\multicolumn{4}{l}{\emph{CDC (38.2M cases, 7 pseudo-sites):}} \\
& Site-level regression & $\beta = +0.55$ & 0.02 \\
& Individual: age per decade & OR $= 2.40$ & $< 0.001$ \\
& Individual: age 80+ & OR $= 9.88$ & $< 0.001$ \\
\midrule
\multicolumn{4}{l}{\emph{Mexico (4.0M cases, 32 states):}} \\
& Site-level regression & $\beta = +1.31$ & $< 0.0001$ \\
& Individual: age 70+ & OR $= 11.26$ & $< 0.001$ \\
& Individual: adjusted & OR $= 7.97$ & $< 0.001$ \\
\bottomrule
\end{tabular}
\end{table}

\begin{figure}[t]
\centering
\includegraphics[width=\textwidth]{figures/fig1_ecological_bias.pdf}
\caption{Site-level ecological regressions for both datasets.
\textbf{(A)}~CDC: 7 race/ethnicity pseudo-sites show a significant
relationship ($\beta = +0.55$, $p = 0.02$), but this linear slope is
incommensurable with the individual-level odds ratio of 9.9.
\textbf{(B)}~Mexico: 32 states show a stronger ecological slope
($\beta = +1.31$, $p < 0.0001$), equally incommensurable with the
individual-level odds ratio of 11.3.  In both cases, ecological
regression reports a linear coefficient where the patient-level truth
is a multiplicative odds ratio an order of magnitude larger.}
\label{fig:ecological}
\end{figure}

\begin{figure}[t]
\centering
\includegraphics[width=0.6\textwidth]{figures/fig2_power.pdf}
\caption{Statistical power of ecological regression depends on the grouping
strategy, not the sample size.  Grouping 105 million CDC records by
race/ethnicity produces pseudo-sites with distinct age distributions,
giving 100\% power.  Random grouping averages out demographic differences:
at $n = 9$ groups, power is 4.8\% (indistinguishable from the
false-positive rate); at $n = 50$, power reaches only 29\%.
Based on 2{,}000 simulations per condition (Supplements~S1 and S6).}
\label{fig:power}
\end{figure}

These results are complementary.  Both datasets show ecological regression
detecting a statistically significant relationship, but in both cases the
ecological coefficient is a linear slope where the patient-level truth is a
multiplicative odds ratio an order of magnitude larger---a different
quantity, not a noisy estimate of the same one.  Any federated analysis
relying on site-level summaries faces this distortion.


\subsection{Why does this happen?}
\label{sec:why}

Two mechanisms drive the distortion.  \emph{Cross-level confounding}: the
7 CDC pseudo-sites differ in sex ratios (15\%--65\% male) and comorbidity
prevalence (0.5\%--14\%), and these confounders are entangled with the
grouping variable at the ecological level.  \emph{Aggregation bias}: the
ecological regression fits a linear slope to group-level proportions,
while the true patient-level relationship is a nonlinear odds ratio.
Noncollapsibility of the odds ratio~\cite{greenland1999noncollapsibility}
compounds both effects---even without confounding, the marginal and
conditional odds ratios differ whenever the outcome is common or the
exposure effect varies across strata.

A simulation (Supplement~S1) isolates the role of grouping strategy:
ecological regression on CDC-matched synthetic data achieves 100\% power
when grouping preserves demographic differences, but only 4.8\% power
under random grouping (Figure~\ref{fig:power})---indistinguishable from
the false-positive rate.

Throughout this paper, ecological $\beta$ (a linear slope), odds ratios
(a multiplicative comparison), and PMIs (observed/expected ratios) each
measure different quantities.

The ecological fallacy arises from the analysis architecture---how data is
aggregated across sites---not from the disease.  COVID-19 serves as a test
case because individual-level records are available for validation, but
the bias applies identically to drug-safety surveillance, federated
genomic association, and any multi-site study that pools site-level
summaries.


\section{Clinical Consequences: The 4CE Neurological Dataset}

The ecological bias demonstrated in Sections~2.1--2.2 has direct clinical
consequences.  The 4CE consortium (Consortium for Clinical Characterization
of COVID-19 by EHR) analyzed neurological outcomes in over 22,000
hospitalized patients across 21 hospitals in 6 countries
\cite{four_ce_neuro2021}, using aggregate site-level data.  The original
4CE study reported site-level proportional morbidity indices and noted
significant between-site heterogeneity; we re-analyze this dataset
using the ecological bias framework from Sections~2.1--2.2 to characterize
\emph{why} the heterogeneity arises and what it implies for the pooled
estimates.  Our critique is not of 4CE specifically---the consortium
acknowledged heterogeneity and used standard methods---but of the broader
practice of pooling site-level summaries when the ecological fallacy applies.


\subsection{CNS involvement and the 350-fold method divergence}

Among well-measured sites ($\text{SE} < 0.5$ on the log scale; $k = 11$--16
per timepoint), the fixed-effects proportional morbidity index (PMI) for CNS
involvement is 1.97 (95\% CI: [1.87, 2.09]) at 30 days.  This estimate must
be interpreted cautiously: $I^2 = 99.8\%$ (Cochran's $Q = 6318$, $p <
10^{-15}$) means essentially all variation is between-site, violating the
common-effect assumption that fixed-effects requires.

The 350-fold divergence between pooling methods (Table~\ref{tab:meta}) is a
diagnostic failure, partly numerical in origin.  The random-effects PMI of
702 is a numerical artifact: several small sites produce PMIs of
$10^{8}$--$10^{18}$ from complete separation (zero cells in contingency
tables), and the DerSimonian--Laird estimator is known to perform poorly
with sparse data and extreme heterogeneity, producing unstable
between-study variance estimates $\hat\tau^2$ that amplify these outliers.  Among the 11 sites with $\text{SE} < 0.4$, the picture is more
consistent (PMIs between 1.4 and 2.8), but even this restricted set shows
$I^2 > 95\%$.  The method divergence signals that this dataset requires
decomposition of heterogeneity rather than a pooled estimate.
Leave-one-site-out analysis confirms this divergence is structural:
dropping each of the 20 sites in turn, the RE/FE ratio ranges from
18 to 984, never falling below 10-fold (Supplement~S8).  No single
site drives the result.

\begin{table}[t]
\centering
\caption{\textbf{The 350-fold divergence between pooling methods is a
diagnostic artifact, not a finding.}  The random-effects PMI of 702 results
from complete separation at small sites (PMIs of $10^{8}$--$10^{18}$)
up-weighted by the DerSimonian--Laird estimator.  No single pooled number
is appropriate for this data ($I^2 = 99.8\%$).  The divergence itself is
the result: it signals that the underlying heterogeneity requires
decomposition, not averaging.}
\label{tab:meta}
\begin{tabular}{llccc}
\toprule
Method & Pooled PMI & 95\% CI & $p$ & $I^2$ \\
\midrule
Fixed-effects (IV) & 1.97 & [1.87, 2.09] & $< 10^{-6}$ & 99.8\% \\
Random-effects (DL) & 702 & [201, 2460] & $< 10^{-6}$ & --- \\
Influence function & 972 & [1.6, 587K] & 0.035 & --- \\
\bottomrule
\end{tabular}
\end{table}

Peripheral nervous system involvement provides a negative control.  PNS
diagnoses show a PMI of 1.04 at 30 days (95\% CI: [0.95, 1.14],
$p = 0.37$), indistinguishable from the null.  If the CNS effect were
driven by sicker patients receiving more diagnostic workups (which would
also detect peripheral neuropathies), PNS diagnoses should show a similar
signal.  They do not.  The specificity to CNS involvement is consistent
with SARS-CoV-2 neuroinvasion through the blood-brain barrier
\cite{song2021neuroinvasion} and central neuroinflammation
\cite{iadecola2020covid}, mechanisms that do not produce peripheral
neuropathy.


\subsection{Age drives between-site variation}

Meta-regression of the CNS severity effect (30-day PMI, log scale) on four
site-level demographics identifies the proportion of elderly patients as the
strongest moderator ($\beta = +14.2$, $\text{SE} = 0.83$,
$p < 10^{-15}$).  The direction is clear: sites serving older populations
show larger neurological effects.  The magnitude should be treated as
directional evidence only.  With ${\sim}20$ sites, Type~M error
\cite{gelman2014typeMS} systematically inflates significant coefficients;
hierarchical shrinkage is the appropriate estimator for effect size, and
we report $\beta = +14.2$ as evidence that elderly proportion is the
dominant moderator, not as an estimate of its strength.

This pattern is consistent with age-dependent blood-brain barrier
vulnerability \cite{montagne2015bbb} and chronic low-grade inflammation
\cite{franceschi2018inflammaging}, but site-level data cannot distinguish
this biological explanation from compositional confounding: sites with more
elderly patients also differ in comorbidity burden, coding practices, and
clinical management.  The E-value for the elderly-proportion moderator
(point estimate: $2.9 \times 10^{6}$; 95\% CI lower bound:
$5.8 \times 10^{5}$) indicates that no plausible unmeasured confounder
could explain the association (Supplement~S7; \cite{vanderweele2017}).

\begin{table}[t]
\centering
\caption{Meta-regression moderators of the CNS severity association (30-day
PMI).  Elderly proportion is the strongest moderator.  Baseline severity
acts in the opposite direction (ceiling effect).  The large $\beta$
magnitudes are likely inflated by Type~M error \cite{gelman2014typeMS} from
$\sim$20 sites.}
\label{tab:moderators}
\begin{tabular}{lcccc}
\toprule
Moderator & $\beta$ & SE & $p$ & Range across sites \\
\midrule
Proportion age $\geq 80$ & $+14.2$ & 0.83 & $< 10^{-15}$ & 0.01--0.43 \\
Proportion female & $+1.13$ & 0.14 & $< 10^{-15}$ & 0.06--0.53 \\
Baseline severity & $-2.91$ & 0.23 & $< 10^{-15}$ & 0.06--0.73 \\
Number of patients & $-0.0003$ & 0.00005 & $< 10^{-9}$ & 32--22{,}789 \\
\bottomrule
\end{tabular}
\end{table}

A site-level interaction between elderly proportion and mortality rate is
significant at 60 and 90 days after FDR correction ($p = 0.024$, $p = 0.021$;
the 30-day value does not survive).  The coefficient magnitudes
($\beta = +32$ to $+38$) are directional evidence only, likely inflated
by Type~M error~\cite{gelman2014typeMS} from $\sim$20 sites.


\subsection{Comorbidity correlations differ systematically across sites}

Elixhauser comorbidity correlation matrices \cite{elixhauser1998} from 20
4CE sites show collective inconsistency when tested with a multivariate
consistency statistic: $z = 25.0$ ($p < 0.0001$), with comparable signals
in CNS-only ($z = 14.4$) and PNS-only ($z = 14.4$) subgroups
(Table~\ref{tab:consistency}).

\begin{table}[t]
\centering
\caption{Cross-site consistency of comorbidity correlation patterns.  The
multivariate consistency statistic detects collective inconsistency
($z = 25$) that pairwise Cochran's $Q$ tests miss entirely (0/1305
significant).}
\label{tab:consistency}
\begin{tabular}{lccc}
\toprule
Test & Statistic & $p$ & Significant pairs \\
\midrule
Consistency (all patients) & $z = 25.0$ & $< 0.0001$ & N/A (collective) \\
Consistency (CNS only) & $z = 14.4$ & $< 0.0001$ & N/A (collective) \\
Consistency (PNS only) & $z = 14.4$ & $< 0.0001$ & N/A (collective) \\
Cochran's $Q$ (per pair) & max $Q = 13.4$ & all $> 0.05$ & 0/1305 \\
\bottomrule
\end{tabular}
\end{table}

No single comorbidity pair varies enough across sites to reach significance
on its own.  The consistency test detects their collective pattern: the
aggregate structure of small per-pair differences is far more extreme than
chance.  Standard meta-analysis assumes the covariate relationships being
adjusted for are constant across sites; this result says they are not.  The most
heterogeneous pairs---HIV-Lymphoma ($Q = 13.4$, $r = 0.09 \pm 0.25$) and
Alcohol-Drugs ($Q = 12.9$, $r = 0.70 \pm 0.24$)---illustrate how coding
and documentation practices drive cross-site inconsistency that the
collective test detects.


\section{Structural Diagnostics Across All Three Datasets}

Two structural diagnostics---the consistency statistic and per-site interaction
analysis---provide partial answers.

\subsection{The consistency test: power requires heterogeneous coverage}

The consistency statistic measures whether per-site correlation matrices
disagree more than expected by chance.  Its power depends on a structural
property of the data: whether different sites report different subsets of
variable pairs (heterogeneous coverage) or all report the same pairs
(homogeneous coverage).

On the CDC data (all 9 race/ethnicity groups including Unknown and
Missing, all reporting all 15 variable pairs), the
permutation null is degenerate: $z = 0$, null standard deviation
$< 10^{-16}$.  On Mexico's 32 states (all reporting all 15 pairs), the
same degeneracy appears: $z = 0.47$, effectively null.
Supplement~A explains this degeneracy and confirms it
empirically on all three datasets.

On the 4CE data, where coverage is heterogeneous (10 sites report
36--296 of 1,305 possible comorbidity pairs; 10 report all 1,305), the
same test produces $z = 25$ ($p < 0.0001$).  The test's power depends on
heterogeneous coverage---a property of the data structure, not the signal.
This is a practical guide for federated networks: the consistency test is
most useful when sites have different measurement capabilities,
which is the norm in multi-national consortia.  A simulation calibration
(Supplement~S10) confirms the test is well-calibrated under the null:
the false-positive rate at $|z| > 2$ is 4.1\% (nominal target
$\leq$ 6\%), and no simulation of 2{,}000 produced $|z| > 25$.


\subsection{Per-site interactions: effect modification varies everywhere}

Per-site logistic regressions reveal genuine variation in the
age $\times$ comorbidity interaction across all three datasets.  In the
CDC data, interaction ORs range 3-fold (0.23--0.69) across 9 groups; in
Mexico, 2-fold (0.19--0.38) across 32 states; in the 4CE data, the
age $\times$ mortality interaction is significant at all three timepoints.
The consistency of this variation across three independent datasets is
direct evidence that pooling loses information about effect modification
(Supplementary Table~S1).


\section{Discussion}

\subsection{Implications for federated analysis infrastructure}

The distortion is worse than noise: a noisy estimate could in principle
be corrected, but a linear slope is not a noisy estimate of an odds
ratio---it is a different quantity entirely.  Aggregation destroys the
within-site covariance structure that individual-level analysis exploits.
Site-level pooling is valid only under restrictive conditions---exchangeable
sites, homogeneous exposure--outcome relationships, no cross-level
confounders---that rarely hold in practice.

The consistency test and per-site interaction analysis offer partial
diagnostics, but neither recovers the individual-level answer.  Three
approaches go further.

\textbf{Federated regression.}
Fed-GLMM~\cite{yan2023fedglmm} and
DataSHIELD~\cite{banerjee2022dssurvival} enable individual-level estimation
while sharing only summary statistics: sites run a standardized script
locally and share coefficients, preserving privacy without aggregation
bias.

\textbf{Large individual-level cohorts.}  N3C~\cite{haendel2021n3c}
(15 million patients, 75+ sites) and ISARIC (705,000 patients, 1,500+
sites) contain the individual records needed for direct patient-level
estimation.

\textbf{Hierarchical models.}  As
Diez-Roux~\cite{diezroux2000multilevel} argued, multilevel models with
site random effects address both the ecological fallacy and Type~M
inflation~\cite{gelman2012multilevel} by shrinking extreme site-level
estimates toward the group mean.

\subsection{Practical guidelines for multi-site studies}

The ecological fallacy is not specific to COVID-19.  Any multi-site
study---genome-wide association, multi-center drug trials, educational
interventions---faces the same risk when the units of analysis (sites)
differ systematically from the units of inference (people).

\paragraph{Report heterogeneity honestly.}
$I^2 = 99.8\%$ means the pooled average is almost meaningless.
Report the full distribution of site-level effects alongside any
pooled estimate.  When $I^2$ exceeds 75\%, treat the pooled value as a
hypothesis, not a conclusion.

\paragraph{Test the structure, not just the average.}
Consistency tests detect coordinated between-site differences that
Cochran's $Q$ (one variable at a time) misses.  Hierarchical models
with site random effects address both the multiplicity problem (through
partial pooling) and Type~M inflation (through shrinkage).  The R
\texttt{metafor} package~\cite{viechtbauer2010metafor} provides these
diagnostics in a single function call (\texttt{rma.mv} with site-level
random effects).

\paragraph{Seek individual-level estimation.}
The ecological fallacy is a problem of aggregation; the definitive
solution is disaggregation.  Modern frameworks make this feasible even
when raw records cannot leave the hospital.
Fed-GLMM~\cite{yan2023fedglmm} achieves near-identical results to
pooled generalized linear mixed models while sharing only summary
statistics: sites run a standardized analysis script locally and share
coefficients (point estimates, standard errors, sample sizes), preserving
privacy while enabling patient-level estimation.
DataSHIELD~\cite{banerjee2022dssurvival} provides privacy-preserving
survival analysis across federated nodes.  These tools should replace
site-level ecological pooling as the default analysis in federated
networks.

\paragraph{Disclose the full analysis menu.}
State how many tests were run and report all results, including nulls.
Pre-register the primary analysis and distinguish it from exploratory
follow-up.


\section{Methods}

\subsection{Reporting guidelines}

This study uses routinely collected health data from three sources.
We follow the RECORD (REporting of studies Conducted using
Observational Routinely-collected health Data) extension of the
STROBE statement where applicable~\cite{benchimol2015record}.
The CDC and Mexico analyses use individual-level data and follow
RECORD items for cohort studies.  The 4CE analysis uses aggregate
site-level summaries and follows STROBE items for cross-sectional
ecological analyses, with deviations noted where site-level data
precludes individual-level reporting.

\subsection{Data sources}

\textbf{CDC Case Surveillance.}  The CDC COVID-19 Case Surveillance Public
Use Dataset contains 38,173,454 confirmed and probable COVID-19 cases reported
to the CDC through 2022.  Each record includes onset date, age group (9
categories: 0--9 through 80+), sex, race/ethnicity (9 categories),
hospitalization, ICU admission, death, and underlying medical condition (binary
flag).  No geographic identifiers are present in the public version.  We used
7 race/ethnicity categories as pseudo-sites, excluding ``Unknown'' and
``Missing'' categories with ambiguous demographic composition ($n$ ranging
from 13,548 to 6,165,036).

\textbf{Mexico Datos Abiertos.}  Mexico's Secretar\'{i}a de Salud publishes
individual-level COVID-19 surveillance data covering all confirmed and
suspected cases nationwide.  We used the January 3, 2022 historical snapshot
(12,425,179 total records, 3,993,464 confirmed COVID-19 cases, 299,581
deaths).  Each record includes age, sex, state of the treating medical unit
(ENTIDAD\_UM, 32 states), death date, and 9 individual comorbidity flags
(diabetes, hypertension, obesity, cardiovascular disease, chronic kidney
disease, COPD, asthma, immunosuppression, smoking).  Cases were classified as
confirmed when CLASIFICACION\_FINAL $\in \{1, 2, 3\}$.  We used state of
medical unit as the site variable, providing 32 real geographic sites with
$n$ ranging from 8,759 (Colima) to 637,408 (Mexico City).

\textbf{4CE Consortium.}  Aggregate site-level data from the 4CE consortium's
Phase 2.1 Neurological Analysis \cite{four_ce_neuro2021}: 21 healthcare sites
across 6 countries (USA, France, Germany, UK, Singapore, Italy), over 22,000
hospitalized COVID-19 patients.  Data included site-level demographics,
Elixhauser comorbidity scores \cite{elixhauser1998}, and proportional morbidity
index (PMI) for CNS and PNS neurological involvement at 30, 60, and 90 days.

\subsection{Statistical analyses}

\textbf{Ecological regression.}  For each dataset, we computed per-site
proportions of elderly patients and mortality rates, then regressed site
mortality on site elderly proportion using ordinary least squares.  For CDC
data, elderly was defined as age 80+; for Mexico, age 70+.

\textbf{Individual-level logistic regression.}  We fit logistic regression
models on the complete individual-level data: (1)~death $\sim$ age;
(2)~death $\sim$ age + sex + comorbidity; (3)~death $\sim$ age $\times$
comorbidity + sex.  For the CDC data, age was coded as the midpoint of the
age group (5, 15, \ldots, 85); for Mexico, age was continuous, dichotomized
at 70.  Odds ratios are reported per decade (CDC) or as binary contrasts
(Mexico).  We used \texttt{statsmodels} v0.14 for estimation.

\textbf{Per-site interaction analysis.}  For each site in each dataset, we
fit a logistic regression with an age $\times$ comorbidity interaction term.
For the 4CE data (no individual records), we fit weighted meta-regressions
with interaction terms on site-level demographics.

\textbf{Multivariate consistency test.}  We computed per-site Pearson
correlation matrices (6 binary variables for CDC and Mexico; 30 Elixhauser
comorbidity categories for 4CE) and measured global consistency as the mean
sum of squared pairwise differences across all site pairs, normalized by
shared coverage.  Significance was assessed by 500 within-dimension
permutations.  The test's power depends on heterogeneous coverage
(Supplement~A).

\textbf{4CE meta-analysis.}  Per-site log-PMI estimates were pooled using
fixed-effects (inverse-variance), DerSimonian--Laird random-effects, and
influence-function methods.  Sites with $\text{SE} \geq 0.5$ were excluded
from primary analyses.  Meta-regression of log-PMI on four site-level
demographics used weighted least squares.  E-values \cite{vanderweele2017}
assess sensitivity to unmeasured confounding.

\subsection{Multiple testing and analysis provenance}

Five primary analyses were pre-specified: ecological and individual-level
regressions on CDC and Mexico data, plus the 4CE fixed-effects PMI.
Twelve secondary 4CE analyses (meta-regression moderators, consistency
tests, interaction tests, PNS negative control, temporal trend) were
developed during exploratory analysis.
Benjamini--Hochberg FDR correction~\cite{benjamini1995fdr} at $q = 0.05$
was applied within the 18-test interaction family; of 3 age--mortality
interaction tests, only the 60- and 90-day values survive.  All null
results are reported (PNS PMI $= 1.04$; consistency test degenerate on
CDC and Mexico; temporal trend $p = 0.16$).  Full multiple testing
accounting appears in Supplement~Table~S1.

Ten supplementary analyses (S1--S10) were pre-registered via a frozen
protocol (commit \texttt{97f7094}) before any results were examined;
two falsification criteria fired (Supplement~B).  No analyses were
conducted and suppressed.

\subsection{Reproducibility}

All code, analysis scripts, and the frozen analysis protocol are
available at
\texttt{github.com/elliottower/ecological-bias-covid}
and are archived on Zenodo (DOI: 10.5281/zenodo.21252372).
The CDC data is publicly downloadable from \texttt{data.cdc.gov} (dataset
\texttt{vbim-akqf}).  The Mexico data is publicly downloadable from
\texttt{datos.gob.mx} (Direcci\'{o}n General de Epidemiolog\'{i}a).


\section{Limitations}

\textbf{Pseudo-sites are not hospitals.}  CDC race/ethnicity pseudo-sites
differ from hospital-based consortium sites: they are not geographically
localized, and the confounders (socioeconomic disparities, healthcare
access) differ from hospital-based confounders (referral patterns, coding
practices, ICU capacity).  Race and ethnicity are social
constructs~\cite{subramanian2009revisiting}; we use CDC categories because
they produce groups with distinct age distributions, the structural property
needed to demonstrate ecological bias.  Mexico's 32 real geographic sites
partially address this limitation; the two datasets together bracket the
range of ecological bias severity.

\textbf{4CE data is aggregate only.}  All 4CE analyses use site-level
summaries.  The age $\times$ mortality interaction coefficients
($\beta = +32$ to $+38$) are likely inflated by Type~M
error~\cite{gelman2014typeMS} from $\sim$20 sites.

\textbf{Consistency tests detect problems, not answers.}  The test flags
unreliable pooling; it cannot recover the individual-level effect.
Its power requires heterogeneous coverage, present in the 4CE data but
absent in the CDC and Mexico data (Supplement~A).

\textbf{Limited variables.}  The CDC dataset has 12 variables; Mexico has
18; hospital datasets such as N3C have hundreds.  Our analysis captures
the direction of the problem but underestimates its magnitude.


\section{Conclusion}

Ecological regression on 42 million COVID-19 patient records from two
countries detects statistically significant age--mortality relationships in
both datasets, but the ecological coefficients (linear slopes of $+0.55$
and $+1.31$) are incommensurable with the individual-level truth
(multiplicative odds ratios of 9.9 and 11.3).  The 4CE consortium's
neurological findings---where meta-analytic methods diverge by 350-fold
and elderly proportion explains most between-site variation---illustrate
the clinical stakes: pooled estimates from this dataset should be treated
as hypotheses, not as established effects.

Structural diagnostics offer partial help.  The consistency test detects
coordinated heterogeneity on the 4CE data ($z = 25$) but is degenerate on
homogeneous-coverage datasets (CDC, Mexico).  Per-site interaction analysis
detects effect modification across all three datasets, confirming that
pooling loses information about how age and comorbidity interact.

The practical recommendation: when heterogeneity is extreme, do not trust the
average.  When individual data is unavailable, test the structure.  When
neither is possible, report site-level findings explicitly as ecological
summaries that may not hold for individual patients---and design the next
study to collect individual-level data.


\appendix

\section{Supplementary Materials}
\label{app:supplementary}

The following supplementary materials are available online:
(A)~Why the consistency test is degenerate under homogeneous
coverage, with empirical confirmation on all three datasets;
(B)~Ten pre-registered supplementary
analyses (S1--S10) with falsification criteria, specified in a frozen
protocol (commit \texttt{97f7094}) before any results were examined---two
falsification criteria fired (S2: Mundlak decomposition, S5: positive
control), both narrowing the CDC demonstration from ``confounding-driven
bias'' to ``small-$k$ uninformativeness'' without affecting the Mexico
or 4CE results; (C)~Power simulation details;
(D)~Duncan--Davis ecological-inference bounds; (E)~Leave-one-site-out
stability analysis; (F)~Consistency-test calibration under the null
(FPR $= 4.1\%$ at $|z| > 2$, well-calibrated).


\section*{Declarations}

\subsection*{Ethics approval and consent to participate}
This study used only publicly available, de-identified datasets.  No
ethics approval or informed consent was required.  The CDC Case
Surveillance Public Use Dataset is released under a public domain license.
The Mexico Datos Abiertos dataset is released by the Secretar\'{i}a de
Salud for public use.  The 4CE consortium data used here consists of
aggregate site-level summaries previously published in
\cite{four_ce_neuro2021}.

\subsection*{Consent for publication}
Not applicable.

\subsection*{Availability of data and materials}
The CDC COVID-19 Case Surveillance Public Use Dataset is available from
\texttt{data.cdc.gov} (dataset identifier \texttt{vbim-akqf}).  The
Mexico COVID-19 surveillance data is available from \texttt{datos.gob.mx}
(Direcci\'{o}n General de Epidemiolog\'{i}a).  The 4CE aggregate data
used in this study is available in the supplementary materials of
\cite{four_ce_neuro2021}.  All analysis code, the frozen
pre-registration protocol, and computed results are available at
\url{https://github.com/elliottower/ecological-bias-covid}.
A permanent archive is deposited at Zenodo
(DOI: \url{https://doi.org/10.5281/zenodo.21252372}).

\subsection*{Competing interests}
The author declares no competing interests.

\subsection*{Funding}
This research received no specific grant from any funding agency in the
public, commercial, or not-for-profit sectors.

\subsection*{Authors' contributions}
ET conceived the study, designed the analyses, wrote the code, performed
all statistical analyses, and wrote the manuscript.  The author read and
approved the final manuscript.

\subsection*{Use of AI-assisted technologies}
Claude (Anthropic) was used as a writing and analysis assistant during
manuscript preparation.  The author directed all analytical decisions,
verified all numerical results against source data, and takes full
responsibility for the content.  The pre-registration protocol,
statistical analyses, and all reported results were independently
validated by the author.  No AI tool is listed as an author.

\subsection*{Acknowledgements}
Not applicable.

\bibliographystyle{unsrtnat}
\bibliography{references}

\end{document}
